\chapter{Conclusiones}

El desarrollo de este trabajo práctico permitió profundizar en el comportamiento y la representación vectorial de las
diferentes conexiones de transformadores trifásicos. A través del simulador interactivo de Aula Moisan y del material
bibliográfico del libro Máquinas Eléctricas de Fraile Mora, se logró comprender cómo las variaciones constructivas y
de conexión eléctrica repercuten directamente en el desfasaje angular de las tensiones secundarias respecto a las primarias.

\section{Análisis de Resultados}
Los resultados obtenidos demuestran que el índice horario no depende únicamente del tipo básico de conexión de los devanados
(estrella, triángulo o zigzag), sino de la polaridad de las conexiones y el conexionado de los terminales homólogos. Se
pueden extraer las siguientes observaciones clave:
\begin{itemize}
    \item La inversión del neutro en una conexión estrella (pasar de Estrella 1 a Estrella 2) o la polaridad del triángulo
    (pasar de Triángulo 1 a Triángulo 2) equivale a un desfase eléctrico de $180^\circ$. Esto se evidenció claramente
    al comparar la conexión Yy0 con la Yy6, o bien al analizar cómo Dy5 difiere de otras conexiones tipo Dy.
    \item La conexión zigzag resulta ser una de las más interesantes desde el punto de vista vectorial, ya que cada fase
    secundaria se compone de la suma vectorial de dos semidevanados bobinados en columnas distintas del núcleo. Las
    configuraciones Zigzag 1 y Zigzag 3 demostraron cómo variando el conexionado interno se pueden obtener desfases
    nulos (Dz0), de $180^\circ$ (Dz6) o bien combinados con la estrella del primario para lograr desfasajes específicos
    como en Yz5 e Yz11.
    \item La utilización del simulador web facilitó la visualización en tiempo real de los triángulos de tensión y de la
    esfera del reloj, validando de manera experimental la teoría analítica descrita en la bibliografía.
\end{itemize}

\section{Cumplimiento de Objetivos}
Todos los objetivos planteados en la introducción del informe fueron plenamente alcanzados:
\begin{itemize}
    \item Se determinó y justificó analíticamente el índice horario para cada una de las nueve configuraciones bajo ensayo.
    \item Se analizó detalladamente la correspondencia entre los esquemas de conexión física de las borneras y sus diagramas
    vectoriales de tensiones.
    \item Se comprendió la relevancia del índice horario en la ingeniería eléctrica práctica. En particular, se reafirmó
    que para posibilitar el acoplamiento en paralelo de dos o más transformadores trifásicos, la igualdad de su índice
    horario es una condición obligatoria y crítica. De no cumplirse esta condición, surgirían diferencias de tensión
    compuesta entre terminales homólogos secundarios que provocarían corrientes de cortocircuito extremadamente elevadas,
    capaces de dañar de forma severa y permanente las máquinas e instalaciones.
\end{itemize}
